\chapter{Naplózás, audit, metrics és hibakeresés}

Egy production webalkalmazás nem attól üzemeltethető, hogy hiba esetén ``kiír valamit a konzolra''. A fejlesztőnek és az üzemeltetőnek ugyanazokat az eseményeket kell tudnia összekapcsolni: melyik HTTP kérés futott, milyen route-ot ért el, milyen státusszal fejeződött be, történt-e biztonsági elutasítás, és maradt-e üzleti szempontból tartós auditnyom.

Az RWLang ehhez négy külön felületet ad:

\begin{itemize}
  \item strukturált \textbf{server log} a technikai és lifecycle eseményekhez;
  \item strukturált \textbf{access log} a HTTP kérésekhez;
  \item strukturált \textbf{audit log} a security és felhasználói aktivitáshoz;
  \item külön listen socketen elérhető \textbf{Prometheus metrics} felületet.
\end{itemize}

Ehhez társul az alkalmazás adatbázisában tárolható, tranzakciós \textbf{business audit trail}. Fontos: az audit log és a business audit nem ugyanaz.

\section{A három naplóosztály szerepe}

Ajánlott production konfiguráció:

\begin{lstlisting}
[observability]
metrics_listen = "127.0.0.1:9090"
allow_public_metrics = false
access_log = true

[logging]
server_file = "/var/log/rwlang/server.log"
access_file = "/var/log/rwlang/access.log"
audit_file = "/var/log/rwlang/audit.log"
stderr = true
\end{lstlisting}

A fájlok célja eltérő.

\begin{description}
  \item[server.log] Indulás, leállás, dependency-hibák, technikai és belső runtime események. Ezt nézd először, ha a szerver nem indul, readiness hibás, vagy egy háttérkomponens problémát jelez.
  \item[access.log] Egy rekord HTTP kérésenként. Ez adja a forgalmi és latency nyomot: method, route, státusz, időtartam, kliens IP, be- és kimenő byte-ok, request ID.
  \item[audit.log] Biztonsági és felhasználói aktivitás: login/logout, auth vagy MFA failure, policy deny, rate-limit deny, resource-profile startup audit és autentikált state-changing action.
\end{description}

Az \code{access_log = false} csak az access logot kapcsolja ki. A security/user audit ettől még megmarad.

\section{Request ID: egy kérés végigkövetése}

Minden sikeresen parse-olt application, static vagy health request szerver-generált azonosítót kap:

\begin{lstlisting}
X-Request-Id: rw-...
\end{lstlisting}

Az access log például ilyen lehet:

\begin{lstlisting}
{"event":"http_request","request_id":"rw-...",
 "method":"GET","route":"articleShow","status":200,
 "duration_ms":7,"client_ip":"203.0.113.5",
 "bytes_in":0,"bytes_out":812}
\end{lstlisting}

A kliens által küldött \code{X-Request-Id} nem authority. Az RWLang saját request ID-t generál, így a támadó nem kényszerítheti rá a szerverre a saját korrelációs azonosítóját.

A hibakeresés alapművelete ezért:

\begin{lstlisting}
# 1. keresd meg a request ID-t az access logban
grep 'rw-...' /var/log/rwlang/access.log

# 2. keresd ugyanazt a technikai logban
grep 'rw-...' /var/log/rwlang/server.log

# 3. security esetben az audit logban is
grep 'rw-...' /var/log/rwlang/audit.log
\end{lstlisting}

A három napló együtt adja a teljes technikai képet.

\section{Mit szabad és mit tilos logolni?}

A log nem debug adatlerakó. Productionban különösen tilos érzékeny request tartalmat automatikusan naplózni.

Az RWLang logfelületére nem kerülhet:

\begin{itemize}
  \item \code{Authorization} header;
  \item Cookie vagy session ID;
  \item CSRF token;
  \item TOTP secret és recovery code;
  \item DB, Redis vagy LDAP credential;
  \item secret-file tartalma;
  \item teljes request body vagy form payload.
\end{itemize}

Ez adatvédelmi és biztonsági invariáns. Ha egy üzleti eseményhez mégis tartósan szükséges információ, ne a request bodyt dumpold: használj explicit, szűk business audit mezőket.

\section{Security audit és HTTP elutasítások}

A trust-boundary elutasításoknál az RWLang külön \code{security_audit} eseményt képez ugyanazzal a request ID-val. Ide tartozik a hibás/untrusted forwarding metadata 400-as elutasítása, a 401/403/421/426/429 security döntések, valamint a CSRF/CORS/origin tiltások finomabb kategorizálása. Így elkülöníthető például:

\begin{itemize}
  \item \code{proxy/invalid_forwarding}: hibás vagy nem trusted forwarding metadata;
  \item \code{auth/unauthorized}: autentikáció hiánya API route-on;
  \item \code{csrf/denied}, \code{cors/denied}, \code{origin/denied}: böngészős trust-boundary tiltás;
  \item \code{policy/forbidden}: route- vagy object-authorization deny;
  \item \code{host/mismatch}, \code{transport/https_required}: Host/TLS boundary sérülése;
  \item \code{rate_limit/denied}: limiter miatti elutasítás.
\end{itemize}

A fejlesztő számára fontos különbség: a 403 nem feltétlen programhiba. Lehet helyes security-policy döntés is. Hibakereséskor mindig nézd meg az audit eseményt és a route policyt is.

\section{Business audit: tartós üzleti változástörténet}

Az access és security log operációs eseménynapló. Egy cikk publikálása, számla jóváhagyása vagy rendelés állapotváltása viszont üzleti állapotváltozás. Ezt az alkalmazás adatbázisában érdemes tranzakciósan rögzíteni.

Példa:

\begin{lstlisting}[language=RWLang]
enum ArticleStatus {
    Review
    Published
}

transaction db {
    Article.publishChanged(tx, id, version)?;
    audit Article id action publish;
        from article.status
        to ArticleStatus.Published
}
\end{lstlisting}

Az audit ugyanabban a DB-tranzakcióban fut, mint a módosítás. Ha az audit rekord nem írható be, a business módosítás is rollbackel. Ez erősebb garancia, mint egy külön log sor kiírása.

Az RWLang a \code{_rw_business_audit} runtime táblába többek között az alábbiakat írja:

\begin{itemize}
  \item runtime-generált eseményazonosító és UTC timestamp;
  \item HTTP request ID;
  \item canonical autentikált actor;
  \item source action;
  \item objektumtípus és objektumazonosító;
  \item statikus üzleti action név;
  \item opcionális előző és új domain állapot.
\end{itemize}

\Needspace{0.30\textheight}
A business audit statement authorizationt nem helyettesít:

\begin{lstlisting}[language=RWLang]
let invoice = Invoice.byId(db, id)?;
authorize invoice owner ownerUsername or role Accountant;

transaction db {
    Invoice.approve(tx, id, version)?;
    audit Invoice id action approve;
        from invoice.status
        to InvoiceStatus.Approved
}
\end{lstlisting}

A helyes sorrend: \textbf{azonosítás -> authorization -> módosítás -> ugyanazon tranzakción belüli audit}.

\section{Prometheus metrics külön management socketen}

A metrics nem application route. Külön listeneren fut:

\begin{lstlisting}
[observability]
metrics_listen = "127.0.0.1:9090"
allow_public_metrics = false
\end{lstlisting}

Lekérdezés ugyanazon a hoston:

\begin{lstlisting}
curl http://127.0.0.1:9090/metrics
\end{lstlisting}

A listener csak a \code{/metrics} GET/HEAD kérést szolgálja ki. Nem fut rajta alkalmazás-route, session vagy auth handler.

Tipikus metrikák:

\begin{lstlisting}
rw_requests_total
rw_responses_5xx_total
rw_auth_failures_total
rw_csrf_failures_total
rw_policy_denials_total
rw_request_timeouts_total
rw_runtime_budget_exceeded_total
rw_rate_limit_denials_total
rw_readiness_failures_total
rw_request_bytes_in_total
rw_response_bytes_out_total
rw_active_connections
rw_log_fallback_total
rw_request_duration_ms_bucket
rw_route_requests_total{route="..."}
rw_route_5xx_total{route="..."}
rw_route_duration_ms_sum{route="..."}
\end{lstlisting}

A \code{rw_policy_denials_total} a security-policy tiltásokat követi: a 403/421/426 döntések mellett a kifejezetten forwarding metadata hibaként felismert 400-as elutasítás is ide számít. A normál malformed request 400 nem security policy denial.

A legfontosabb szabály a \textbf{bounded cardinality}. A \code{route} label compiler-owned route név vagy fix reserved név. Raw URL, path paraméter, IP, username, tenant, session ID és request ID soha ne legyen metric label. Ellenkező esetben egy támadó korlátlan idősor-számot hozhatna létre.

\Needspace{0.42\textheight}
\section{Apache vagy Nginx mellett ki scrape-elje a metricset?}

A publikus webes proxy és a metrics listener külön trust boundary. Tipikus elrendezés:

\begin{lstlisting}
Internet
   |
   v
Apache/Nginx :443
   |
   v
127.0.0.1:8080  RWLang server

Prometheus agent
   |
   v
127.0.0.1:9090  RWLang metrics
\end{lstlisting}

A metrics portot ne proxyd ki automatikusan ugyanarra a publikus domainre. Ha távoli Prometheus szervernek kell elérnie, használj privát management hálózatot, ACL-t vagy külön autentikált monitoring proxyt.

\section{Health: liveness és readiness nem ugyanaz}

Az RWLang beépített szerver endpointjai:

\begin{lstlisting}
GET  /health/live
HEAD /health/live

GET  /health/ready
HEAD /health/ready
\end{lstlisting}

A liveness csak azt jelzi, hogy a szerver event loop él. Nem hív DB-t, Redist vagy LDAP-ot.

A readiness azt mutatja, hogy a kiszolgáláshoz szükséges konfigurált runtime dependency-k elérhetők. DB vagy Redis hiba/timeout esetén \code{503 Service Unavailable} lehet a válasz.

Ezért load balancernél:

\begin{itemize}
  \item a \textbf{liveness} alapján ne vedd ki rögtön az instancet egy rövid DB kiesés miatt;
  \item a \textbf{readiness} alapján viszont ideiglenesen kivehető a forgalomból.
\end{itemize}

A health endpoint nem üzleti monitoring API: szándékosan minimális státuszt közöl.

\Needspace{0.46\textheight}
\section{Logrotate és SIGHUP}

Az RWLang file logging append módban ír. Productionban a logrotációt az operációs rendszerre bízd. A repository példája: \code{examples/logrotate/rwlang}.

Egyszerűsített minta:

\begin{lstlisting}
/var/log/rwlang/server.log /var/log/rwlang/access.log {
    daily
    rotate 14
    compress
    delaycompress
    missingok
    notifempty
    create 0640 rwlang rwlang
    sharedscripts
    postrotate
        systemctl kill -s HUP --kill-who=main rwlang.service || true
    endscript
}
\end{lstlisting}

A \code{SIGHUP} jelentése tudatosan korlátozott. Behind-proxy módban újranyitja a logfájlokat és tranzakciósan újraépíti a domain/application hosting állapotot; nem általános process-szintű \code{server.toml} rekonfiguráció. Az alkalmazásforrás változását normál esetben az automatikus source-reload supervisor kezeli.

Konfigurációmódosításnál:

\begin{lstlisting}
/usr/local/bin/rwlang-server \
  --config /usr/local/etc/rwlang/server.toml \
  --check-config

systemctl restart rwlang.service
\end{lstlisting}

Ez auditálhatóbb és kiszámíthatóbb, mint egy részlegesen sikerülő hot reload.

\section{Mi történik, ha a logfájl nem írható?}

A file logging bounded writer queue-t használ. Ha a queue megtelik, a file sink kiesik vagy a writer nem érhető el, az RWLang stderr fallbacket használ és növeli a következő metrikát:

\begin{lstlisting}
rw_log_fallback_total
\end{lstlisting}

Ezt érdemes riasztási feltételként kezelni. Ha növekszik, lehet:

\begin{itemize}
  \item jogosultság- vagy filesystem-hiba;
  \item megtelt lemez;
  \item hibás logrotate/reopen;
  \item tartós log-processing nyomás.
\end{itemize}

A logfájl útvonalának parent directory-ja létezzen. Unix rendszeren a runtime a logfájlt \code{0640} móddal, symlink-követés nélkül nyitja meg; az operatornak ettől még helyesen kell kialakítania a könyvtárjogosultságokat.

\section{Gyakorlati production hibakeresési folyamat}

Egy 5xx vagy ``lassú az oldal'' bejelentésnél ne találgass. Használj reprodukálható sorrendet.

\begin{enumerate}
  \item \textbf{Liveness:} fut-e maga a process?
  \item \textbf{Readiness:} elérhető-e DB/Redis?
  \item \textbf{Access log:} melyik route, státusz, duration és request ID érintett?
  \item \textbf{Server log:} ugyanahhoz a request ID-hoz tartozik-e dependency/runtime hiba?
  \item \textbf{Audit log:} security policy vagy rate limit utasította-e el?
  \item \textbf{Metrics:} egyedi eset vagy trendszerű 5xx/latency növekedés?
  \item \textbf{Reverse proxy:} a hiba az Apache/Nginx előtt, benne vagy az RWLang mögött keletkezik?
\end{enumerate}

Példa gyors ellenőrzésre:

\begin{lstlisting}
curl -i https://www.example.com/health/live
curl -i https://www.example.com/health/ready
# Proxy deploymentnel plusz belso upstream-proba:
curl -i http://127.0.0.1:8080/health/ready
curl -s http://127.0.0.1:9090/metrics | grep rw_responses_5xx_total
journalctl -u rwlang.service --since "10 min ago"
\end{lstlisting}

Ha a loopback RWLang endpoint jó, de a publikus domain rossz, a következő vizsgálati réteg az Apache/Nginx, TLS vagy forwarded-header policy. Ha a loopback readiness is 503, előbb a backend dependency-ket vizsgáld.

\section{Minimum monitoring baseline}

Egy kisebb production telepítésnél is legyen legalább riasztás az alábbiakra:

\begin{itemize}
  \item readiness tartósan nem ready;
  \item 5xx arány vagy számláló hirtelen növekedése;
  \item request latency érdemi romlása;
  \item auth/policy/rate-limit denial szokatlan kiugrása;
  \item \code{rw_runtime_budget_exceeded_total} növekedése;
  \item \code{rw_log_fallback_total} növekedése;
  \item process restart loop vagy váratlan shutdown.
\end{itemize}

A konkrét thresholdot nem a nyelvnek kell hardcode-olnia. A forgalom, SLO és üzleti kritikalitás alapján az üzemeltető állítja be.

\section{Retention, SIEM és auditvédelem}

A retention szervezeti és jogi policy. Az RWLang nem kényszerít rá univerzális megőrzési időt.

A helyi \code{audit.log} fontos operációs nyom, de önmagában nem erős tamper-evident archívum, mert a service host adminisztrátora vagy megfelelő filesystem jogosultságú szereplő módosíthatja. Magasabb assurance esetén továbbítsd centralizált, hozzáférés-védett vagy append-only log/SIEM rendszerbe.

A business audit adatbázis-tábla ugyancsak nem véd a DB-admin közvetlen módosítása ellen. A két réteg együtt ad jó alkalmazási nyomot, de compliance környezetben a külső, független logarchívum marad fontos.

\section{Ellenőrzőlista}

Production indulás előtt ellenőrizd:

\begin{itemize}
  \item külön server/access/audit log útvonal van;
  \item a logkönyvtár jogosultságai szűkek;
  \item request body és secret nincs logolva;
  \item a metrics listener nem publikus véletlenül;
  \item a Prometheus labelk bounded cardinality-jűek;
  \item a liveness és readiness külön van kezelve;
  \item logrotate után SIGHUP történik, configmódosítás után restart;
  \item \code{rw_log_fallback_total} monitorozva van;
  \item üzleti workflow esetén a business audit ugyanabban a DB-tranzakcióban fut;
  \item az audit log központi továbbításáról és retentionjéről van dokumentált döntés.
\end{itemize}
